%==================================================================== 
% This file is style for thesis.
% Modified by  maeda@taka.is.uec.ac.jp ' 2012/01/13
%==================================================================== 
%-------------------------------------------------------------------- 
%\documentclass[12pt,epsbox,subfigure]{jbook}
\documentclass[12pt,epsbox,epsf,eclepsf,subfigure,oneside,dvipdfmx]{jbook}
%
% Style of page
%
% header and footer
% \pagestyle{headings}
% \pagenumbering{arabic}

\pagestyle{plain}
\usepackage[centrefoot]{pageno}

\markright{}
%
% A4 paper:  W=211mm, H=297mm.
\setlength{\textheight}{230mm}
\setlength{\topmargin}{-10mm}
%\setlength{\headheight}{10mm}
\setlength{\headsep}{15mm}
%\setlength{\footheight}{0mm}
\setlength{\footskip}{10mm}
\setlength{\textwidth}{165mm}
\setlength{\oddsidemargin}{0mm}
\setlength{\evensidemargin}{0mm}
\setlength{\parindent}{10pt}
\setlength{\parskip}{0.75pt}
\setlength{\unitlength}{1.5mm}
\setlength{\unitlength}{1mm}
%-------------------------------------------------------------------- 
%-------------------------------------------------------------------- 
\renewcommand{\textfraction}{0.2}
\renewcommand{\topfraction}{0.8}
\renewcommand{\bottomfraction}{0.8}
\newtheorem{lemma}{Lemma}[section]
\newtheorem{proposition}{Proposition}
\newtheorem{corollary}{Corollary}[section]

% for mathematical symbols.
\newcommand{\dgr}{^\circ}
\newcommand{\atan}{\mbox{atan}}
\newcommand{\atant}{\mbox{atan2}}
\newcommand{\mbold}[1]{\mbox{\boldmath$#1$}}
\newcommand{\hatbold}[1]{\hat{\mbox{\boldmath$#1$}}}

% for minipage environment.
\newlength{\minitwocolumn}
\setlength{\minitwocolumn}{0.5\textwidth}
\addtolength{\minitwocolumn}{-0.5\columnsep}

% shee definition.
\def\A4sheet{\width=210mm \height=297mm}
\def\A5sheet{\width=148mm \height=210mm}
\def\B4sheet{\width=257mm \height=363mm}
\def\B5sheet{\width=182mm \height=257mm}
\def\bibname{参考文献}  %article�Atarticle�N���X
%-------------------------------------------------------------------- 
\renewcommand{\baselinestretch}{1.5}
% \usepackage[draft]{graphicx}
\usepackage{graphicx}
\usepackage{subcaption}
\usepackage{float}
\usepackage{mathrsfs} 
\usepackage{amsmath}
\usepackage{amsfonts}
%\usepackage{slashbox}
%\usepackage{mathpazo}
%\usepackage{txfonts}

% algorithm
\usepackage{algorithm}
\usepackage{algpseudocode}

\usepackage{tabularx} % For tables with auto-wrapping columns
% Add because main language is not English
\usepackage{enumitem}
\setlist[itemize]{label=-}

\usepackage{bm}
\usepackage{cite}
\usepackage{url}
\usepackage{xcolor}

%\bibliographystyle{jplain}
% \usepackage{fancyhdr}

%%%%%%%%%%%%%%%%%%%%
\usepackage{url} 
\usepackage{boites} 
\usepackage{colortbl}
\usepackage{multirow}
\usepackage{listings,jlisting}


% %%%%Table
\usepackage{booktabs}
\usepackage{tabularx}
\usepackage{array}
\usepackage{makecell}


\definecolor{orange}{rgb}{1.0, 0.65, 0.0}

\newcommand{\orange}[1]{\multicolumn{1}{>{\columncolor{orange}}c|}{#1}}
%%%%%%%%%%%%%%%%%%%%

% \usepackage{subfigure}
\newcommand{\csubfloat}[2][]{%
  \makebox[0pt]{\subfloat[#1]{#2}}%
}
\newcommand{\centerhfill}[1][\quad]{\hspace{\stretch{0.5}}#1\hspace{\stretch{0.5}}}

% for Table
\usepackage{booktabs}

\def\MARU#1{{\ooalign{\hfil#1\/\hfil\crcr\raise.167ex\hbox{\mathhexbox20D}}}} 
\renewcommand{\labelitemi}{�E}
% \newcommand\subcaption[1]{\begin{center}#1\end{center}}

\newcommand{\namefingerpoint}{指先位置}
\newcommand{\namefingerforce}{指先力}
\newcommand{\nameholdingpointset}{把持位置}
\newcommand{\namesearchspace}{探索空間}
\newcommand{\namegraspforce}{把持力} % 指先力の集まり
\newcommand{\namefingerpointENG}{finger position}
\newcommand{\namefingerforceENG}{finger force}
\newcommand{\nameholdingpointsetENG}{holding point set}
\newcommand{\namesearchspaceENG}{search space}
\newcommand{\namegraspforceENG}{holding force}

\newcommand{\namegraspforcetitleENG}{HoldingForce}
\newcommand{\namefingerpointcapENG}{Finger position}
\newcommand{\namefingerforcecapENG}{Finger force}
\newcommand{\nameholdingpointsetcapENG}{Holding point set}
\newcommand{\namesearchspacecapENG}{Search space}
\newcommand{\namegraspforcecapENG}{Holding force} % TODO: 重大问题，原文中有2种解释．．．6*1的合力 以及 3n*1指力列向量;

\usepackage{xparse}
% Revision marks for the changes made after the preliminary review.
% Use \showthesisrevisionsfalse to produce the clean final manuscript.
\newif\ifshowthesisrevisions
\showthesisrevisionstrue
\definecolor{thesisrevision}{RGB}{200,0,0}
\NewDocumentCommand{\revise}{+m}{%
  \ifshowthesisrevisions
    {\color{thesisrevision}#1}%
  \else
    #1%
  \fi
}
\newenvironment{revisionblock}
  {\ifshowthesisrevisions\color{thesisrevision}\fi}
  {}

\NewDocumentCommand{\equfingerpoint}{ O{} O{} }{\bm{p}_{\mathrm{#1}#2}}
\NewDocumentCommand{\equholdingpointset}{ O{} O{} }{{P}_{\mathrm{#1}#2}}  % 把持位置
\newcommand{\equsearchspace}{\mathbb{P}}
\NewDocumentCommand{\equfingerforce}{ O{} }{\bm{f}_{\mathrm{f}#1}}
\NewDocumentCommand{\equgraspforcevector}{ O{} }{\bm{f}_{\mathrm{#1}}}
\newcommand{\equgraspforce}{\bm{f}}
\newcommand{\equgraspforcelist}{\mathbb{F}}

\NewDocumentCommand{\fix}{ O{} }{\textcolor{black}{#1}}
\NewDocumentCommand{\todofix}{ O{} }{\textcolor{black}{#1}}



